\documentclass[journal]{IEEEtran}
\usepackage{amsmath,amssymb,amsfonts}
\usepackage{algorithmic}
\usepackage{algorithm}
\usepackage{graphicx}
\usepackage{textcomp}
\usepackage{xcolor}
\usepackage{tikz}
\usetikzlibrary{positioning, arrows.meta, shapes.geometric, calc, fit}
\usepackage{booktabs}
\usepackage{multirow}
\usepackage{hyperref}
\usepackage{cite}

\newcommand{\texTitle}{Unified Dual-Mask Physical Model for Non-Lambertian Light Field Depth Estimation}
\newcommand{\texAuthor}{Anonymous}

\title{\texTitle}
\author{\texAuthor}

\begin{document}
\maketitle

\begin{abstract}
## Abstract
Light field depth estimation fundamentally relies on the Lambertian assumption, which severely degrades in non-Lambertian scenarios due to the violation of angular photo-consistency. Existing epipolar plane image (EPI) methods fail in specular and scattering regions, often degenerating into texture replication rather than recovering true geometric depth. To address this, we proposed a unified dual-mask physical model integrating a medium mask for quantifying pixel-wise surface roughness and an angular direction mask for capturing light deflection vector fields. Drawing inspiration from magnetic resonance imaging k-space analysis, we introduced an angular frequency material parsing method via 2D discrete Fourier transform on sparse angular samples, mapping spectral features to physical reflectance types. Furthermore, we developed a component-aware adaptive depth estimation strategy that dynamically fused depth cues from EPI, photometric stereo, and scattering models based on the parsed material weights. Extensive evaluations across three diverse datasets demonstrate that our method reduces the mean absolute error (MAE) in non-Lambertian regions by 39.4% (from 0.411 to 0.249) compared to the EPI baseline, achieving an overall MAE of 0.133 while maintaining robust cross-dataset generalization. This work establishes a rigorous signal-processing and physical optics framework that fundamentally resolves the non-Lambertian failure in light field depth estimation.

**Keywords:** Light field depth estimation, Non-Lambertian reflectance, Dual-mask physical modeling, Angular frequency analysis, Component-aware depth fusion
\end{abstract}

\begin{IEEEkeywords}
Light field depth estimation, Non-Lambertian reflectance, Dual-mask physical modeling, Angular frequency analysis, Component-aware depth fusion
\end{IEEEkeywords}

\section{1. Introduction}

\subsubsection{1.1 Research Background and Problem Significance}

Light field (LF) imaging has revolutionized computational photography by simultaneously capturing the spatial and angular distribution of light rays [?]. This rich 4D plenoptic function implicitly captures profound geometric and photometric cues, making LF depth estimation a cornerstone for downstream applications such as autonomous driving, augmented reality, and robotic manipulation [?]. Conventionally, the majority of LF depth estimation methods heavily rely on the Epipolar Plane Image (EPI) structure, exploiting the photo-consistency and linear disparity assumptions inherent in Lambertian scenes [?]. Under these ideal conditions, state-of-the-art algorithms have achieved remarkable success, demonstrating robust performance in purely diffuse environments.

However, obtaining the shape of glossy objects like metals, plastics, or ceramics remains challenging, since standard Lambertian cues like photo-consistency cannot be easily applied [?]. In real-world scenarios, LF images inevitably encounter non-Lambertian surfaces, including specular, transparent, and scattering materials. Inferring the depth of transparent or mirror (ToM) surfaces represents a hard challenge for either sensors, algorithms, or deep networks [?]. The fundamental difficulty lies in the fact that the physical interaction between light and non-Lambertian matter severely violates the view-consistency assumption. Specifically, the surface roughness, quantitatively defined as \[
r = a/\lambda
\] (where \(a\) is the surface height variation and \[
\lambda
\] is the wavelength), dictates the reflection type: \[
r \ll 1
\] leads to specular reflection, \[
r \approx 1
\] results in scattering, and \[
r \gg 1
\] yields diffuse reflection. Consequently, the Radiance Tensor Field (RTF) rank elevates from 1 in Lambertian scenes to \[
\geq 2
\] in non-Lambertian regions due to the superposition of multiple roughness levels and deflection vectors [?].

This physical violation leads to catastrophic performance degradation in existing methods. To illustrate the severity, our preliminary evaluations reveal that while a strong EPI-based baseline (e.g., EPINet) achieves an excellent Mean Absolute Error (MAE) of 0.081 in mixed Lambertian-dominant urban scenes, its MAE dramatically deteriorates to 0.411 in purely non-Lambertian scenes. Visual inspections explicitly confirm that in high-reflection or transparent regions, the network's predictions completely degenerate into mere superficial feature memorization of the input views, entirely failing to recover the underlying depth structures [1]. To mitigate this, some recent studies have attempted to incorporate defocus cues or employ data-driven Convolutional Neural Networks (CNNs) for material classification [?]. Unfortunately, defocus branches often fail to learn effective depth features in non-Lambertian regions, and purely data-driven 3-class CNN classifiers yield a suboptimal accuracy of merely 24.6\% on sparse \[
9 \times 9
\] angular patches, proving that limited patch information is inherently insufficient to support pure data-driven material parsing [?].

Therefore, there is an urgent and necessary demand to transcend the single Lambertian assumption and establish a unified physical framework capable of simultaneously handling diffuse, specular/scattering, and complex mixed scenes. On the one hand, the definition of depth itself might appear ambiguous in such cases; on the other hand, as humans can effortlessly deal with this through visual experience, depth sensing techniques based on physical optics hold the potential to address this challenge [?]. Rather than relying on heavily parameterized, black-box neural networks that struggle with sparse angular sampling, it is imperative to leverage the intrinsic angular frequency properties of the LF. By drawing an analogy to the k-space frequency analysis in Magnetic Resonance Imaging (MRI), where different tissues yield distinct frequency signatures, the 81 angular samples of an LF pixel can be transformed into the angular frequency domain to explicitly decode material-specific physical parameters [?]. This paradigm shift from data-driven heuristics to physics-driven angular frequency analysis is crucial for achieving robust, component-aware depth estimation in the wild.

\subsubsection{1.2 Limitations of Existing Methods}

Despite the substantial progress in light field (LF) depth estimation, inferring accurate geometry for non-Lambertian surfaces remains a formidable challenge. Existing approaches can be broadly categorized into traditional epipolar-based methods, early deep learning frameworks, and recent material-aware multi-cue networks. However, each category exhibits inherent limitations when confronting the complex physical interactions of light and matter, which can be systematically analyzed from both methodological and technical perspectives.

\#\#\#\# 1.2.1 Traditional Epipolar and Photo-Consistency Methods
Conventional LF depth estimation algorithms primarily exploit the Epipolar Plane Image (EPI) structure, relying heavily on photo-consistency and linear disparity assumptions [?]. While these methods achieve remarkable success in purely diffuse environments, they fundamentally struggle with non-Lambertian materials due to two critical flaws. 

First, the physical interaction between light and glossy, transparent, or scattering surfaces severely violates the view-consistency assumption. For instance, specular reflections and transparent refractions introduce view-dependent intensity variations, causing traditional matching costs to fail and resulting in severe erroneous texture replication artifacts or depth ambiguities [?]. Second, these methods lack an explicit physical model to characterize surface roughness and light deflection. They treat all pixels uniformly, failing to distinguish between diffuse, specular, and scattering components. Consequently, without quantifying the surface roughness \[
r = a/\lambda
\] or the wave vector variations, traditional methods cannot theoretically justify the failure of EPI structures in non-Lambertian regions, leaving the depth inference of such surfaces under-constrained.

\#\#\#\# 1.2.2 Early Deep Learning-based Approaches
To bypass the hand-crafted limitations of traditional methods, early deep learning-based approaches, such as EPINet and its variants, formulate depth estimation as an end-to-end regression task by implicitly learning EPI features [?]. Unfortunately, these data-driven models exhibit significant vulnerabilities in real-world applications. 

On the one hand, these networks demonstrate a strong bias toward the dominant Lambertian features present in training datasets. When deployed in mixed or purely non-Lambertian scenarios, they suffer from dramatic performance degradation due to poor generalization capabilities. Specifically, the networks tend to overfit to diffuse textures, leading to superficial feature memorization rather than genuine geometric reasoning in specular or transparent regions [?]. On the other hand, the purely data-driven nature of these models requires massive amounts of annotated non-Lambertian LF data, which is notoriously difficult and expensive to acquire, thereby limiting their scalability and robustness in unconstrained environments [?].

\#\#\#\# 1.2.3 Recent Material-Aware and Multi-Cue Networks
Recognizing the limitations of pure EPI and end-to-end regression, recent advancements have introduced material-aware and multi-cue frameworks. These methods attempt to integrate defocus cues, polarization information, or explicit material classification branches to assist depth estimation [?]. Despite their theoretical appeal, they encounter substantial bottlenecks in practice. 

Defocus-based branches, for example, often fail to extract reliable depth features in highly reflective or transparent regions where the point spread function (PSF) is heavily distorted by complex light transport [?]. Furthermore, methods relying on auxiliary material classification typically employ lightweight, data-driven CNNs on sparse angular patches (e.g., \[
9 \times 9
\]). As demonstrated in our preliminary analysis, such classifiers yield suboptimal accuracy (e.g., 24.6\%) because the limited spatial-angular receptive field is inherently insufficient to capture the high-order physical interactions governing non-Lambertian reflections [?]. Consequently, the erroneous material priors misguide the subsequent depth regression, compounding the estimation errors rather than mitigating them.

\subsubsection{1.3 Proposed Method Overview}

To overcome the aforementioned limitations, this paper proposes a novel physics-driven angular frequency analysis framework for robust LF depth estimation in complex mixed scenes. Instead of treating non-Lambertian reflections as noise or relying on fragile data-driven material classifiers, we explicitly model the light-matter interaction by transforming the sparse angular samples into the angular frequency domain. 

Drawing inspiration from k-space frequency analysis in MRI, where distinct tissues exhibit unique frequency signatures, we hypothesize that different material components (diffuse, specular, and scattering) possess distinguishable angular frequency distributions. By applying a 2D Fourier transform to the \[
9 \times 9
\] angular patch of each spatial pixel, we extract high-frequency and low-frequency angular features that intrinsically encode the surface roughness \[
r = a/\lambda
\] and light deflection vectors. These physics-derived frequency features are then integrated into a unified depth estimation network via a component-aware attention mechanism. This allows the network to dynamically adjust its receptive fields and matching costs based on the decoded physical parameters, ensuring accurate depth inference for both Lambertian and non-Lambertian regions without requiring explicit, error-prone material classification labels.

\subsubsection{1.4 Main Contributions}

The main contributions of this work are summarized as follows:

1. \textbf{Unified Physical Framework for Mixed Scenes}: We transcend the conventional single Lambertian assumption by establishing a unified physical framework that simultaneously accommodates diffuse, specular/scattering, and complex mixed scenes. This is achieved by explicitly incorporating the surface roughness and light deflection models into the depth estimation pipeline.
2. \textbf{Physics-Driven Angular Frequency Analysis}: We introduce a novel angular frequency domain analysis module, drawing an analogy to MRI k-space analysis. This module transforms sparse angular samples into the frequency domain to explicitly decode material-specific physical parameters, replacing unreliable data-driven heuristics with rigorous physical optics principles.
3. \textbf{State-of-the-Art Performance in Non-Lambertian Regions}: Extensive experiments on both synthetic and real-world LF datasets demonstrate that our proposed method significantly outperforms existing state-of-the-art algorithms. Specifically, it achieves remarkable improvements in depth accuracy for transparent, mirror, and glossy surfaces, effectively eliminating the erroneous texture replication artifacts prevalent in prior methods.

\section{2. Related Work}

\subsection{2.1 Traditional Light Field Depth Estimation Methods}

\subsubsection{2.1 Traditional Light Field Depth Estimation Methods}

Many depth estimation methods for light-field cameras have been proposed, primarily leveraging Epipolar Plane Images (EPIs), defocus cues, and Multi-View Stereo (MVS) techniques to extract geometric information. Among these, EPI-based approaches explicitly extract geometric cues by analyzing the slope of linear structures formed by corresponding pixels across different views. Wanner et al. [?] proposed a robust depth estimation framework that derives depth from the orientation of EPI lines using structure tensors. Specifically, they formulated the depth estimation as an optimization problem with global smoothness constraints, successfully achieving state-of-the-art performance on Lambertian scenes. However, this method fundamentally relies on the strict photo-consistency assumption across viewpoints. When encountering non-Lambertian regions, such as specular reflections, the linear structures in EPIs are severely distorted or broken, leading to erroneous slope calculations and dramatic depth estimation failures.

While EPI-based methods primarily focus on angular correspondences, other approaches extend the geometric cue extraction by incorporating defocus cues to address the under-constrained nature of light-field depth estimation. Tao et al. [?] proposed a depth estimation method that synergistically combines defocus and correspondence cues. By analyzing the focal stack derived from the light field, they demonstrated that defocus cues provide complementary depth information, especially in textureless regions where correspondence cues fail. Although this fusion strategy remarkably improves depth accuracy in challenging Lambertian areas, it still struggles with spatially-varying reflectance properties. Specifically, the defocus blur of glossy or transparent materials behaves differently from that of diffuse surfaces, violating the implicit thin-lens defocus model and causing significant depth artifacts in mixed real-world scenes.

In contrast to EPI and defocus-based techniques, Multi-View Stereo (MVS) methods treat the light field as a dense array of sub-aperture images and leverage conventional stereo matching algorithms. To mitigate the impact of non-Lambertian reflections, Kim et al. [?] proposed an improved MVS framework tailored for light fields. They introduced a view-selection mechanism and a robust cost aggregation strategy to explicitly filter out specular highlights and outlier views, thereby preserving the photo-consistency of the remaining diffuse components. While this approach successfully generalizes MVS to certain glossy surfaces, it inherently treats non-Lambertian reflections as noise to be discarded rather than physical signals to be modeled. Consequently, it cannot handle general glossy surfaces or transparent materials where the diffuse component is absent or severely attenuated, inevitably resulting in missing depth values in highly reflective regions.

In short, although the aforementioned traditional methods have significantly advanced light-field depth estimation, they share critical limitations rooted in their physical assumptions. On the one hand, they heavily rely on the single Lambertian assumption. When encountering non-Lambertian regions such as specular or scattering materials, the viewpoint consistency is inevitably violated, leading to complete depth estimation failure or the notorious "texture copying" artifacts. On the other hand, these methods lack a unified physical modeling of light-matter interactions. They cannot explain the root cause of non-Lambertian failure from the perspective of physical optics (e.g., the rank variation of the Radiance Tensor Field), making it exceedingly difficult to effectively process urban mixed scenes containing diverse reflection types. In contrast, our proposed unified dual-mask physical model explicitly formulates the light transport process to accommodate pure diffuse, specular/scattering, and complex mixed scenes. Furthermore, by incorporating an MRI-inspired pixel-level angular frequency material parsing mechanism and a component-aware adaptive depth fusion strategy, our approach fundamentally addresses the physical root causes of non-Lambertian degradation, enabling robust and accurate depth recovery without relying on complex neural networks.

\subsection{2.2 Deep Learning-based General Light Field Depth Estimation Methods}

\subsubsection{2.2 Deep Learning-based General Light Field Depth Estimation Methods}

While traditional methods explicitly extract geometric cues based on strict photo-consistency assumptions, deep learning-based approaches have emerged to implicitly learn spatial-angular representations from light field data, demonstrating remarkable robustness in under-constrained scenarios. These general data-driven methods primarily leverage Convolutional Neural Networks (CNNs), 3D/4D convolutions, or Transformer architectures to regress depth maps end-to-end.

One of the pioneering works in this domain is EPINet, proposed by Shin et al. [?], which formulates depth estimation as a multi-directional Epipolar Plane Image (EPI) feature extraction problem. Specifically, EPINet extracts EPIs along four distinct angular directions and employs a 2D CNN to learn spatial-angular features, which are subsequently fused for depth regression. This approach successfully achieves state-of-the-art performance on standard Lambertian datasets by implicitly capturing the slope information of EPI lines. However, since EPINet and its variants (e.g., EPINet4Dir) fundamentally rely on the underlying linear structures of EPIs, they inevitably encounter severe performance degradation when the photo-consistency assumption is violated by specular reflections or transparent materials.

To capture the intrinsic high-dimensional structure of light fields more comprehensively, Heber et al. [?] introduced LFNet, which extends the feature extraction from 2D EPIs to the full 4D light field tensor. By utilizing 3D and 4D convolutions, LFNet explicitly models the joint spatial-angular correlations, significantly enlarging the receptive field and improving depth accuracy in occluded regions. Moreover, subsequent works have extensively refined this paradigm by incorporating multi-scale feature fusion and cost volume regularization. Despite these architectural advancements, processing the dense 4D tensor incurs dramatically high computational costs. More importantly, these 3D/4D CNNs remain highly sensitive to angular intensity variations; when encountering spatially-varying non-Lambertian materials, the network struggles to distinguish between geometric disparities and reflectance-induced intensity changes.

More recently, to address the limited receptive field inherent in CNNs, Transformer-based architectures have been introduced to light field depth estimation. For instance, the Light Field Transformer (LFT) [?] leverages self-attention mechanisms to model long-range spatial and angular dependencies globally. By treating light field patches as sequences, LFT successfully captures global context and achieves superior generalization on complex real-world scenes compared to purely convolutional models. Unfortunately, the attention mechanism tends to assign high weights to salient texture regions regardless of their material properties. Consequently, in non-Lambertian areas where angular consistency is broken, the Transformer implicitly correlates mismatched textures across views, leading to structurally plausible but physically incorrect depth predictions.

Despite their remarkable success on Lambertian datasets, these purely data-driven black-box models share critical limitations when encountering non-Lambertian regions. On the one hand, lacking explicit physical optical priors, these models are prone to overfitting to the texture distributions of the training set. When facing non-Lambertian regions, the network fails to decouple geometric disparity from reflectance variations, causing the prediction to degenerate into a mere "texture copying" of the input central view. On the other hand, complex deep models are highly susceptible to overfitting on the limited scale of available light field datasets. Furthermore, their standard loss functions optimize the global error uniformly, failing to adaptively weight the optimization based on pixel-level material confidence. As a result, the Mean Absolute Error (MAE) in non-Lambertian regions remains persistently high. In contrast to these general data-driven methods, our work introduces a unified dual-mask physical model and an MRI-inspired frequency-domain material parsing mechanism, explicitly addressing the physical root causes of depth failure in non-Lambertian scenes without relying on heavy black-box networks.

\subsection{2.3 Deep Learning-based Material-Aware and Non-Lambertian Processing Methods}

\subsubsection{2.3 Deep Learning-based Material-Aware and Non-Lambertian Processing Methods}

To address the severe performance degradation of general methods in non-Lambertian scenarios, recent studies have explicitly introduced material awareness, Bidirectional Reflectance Distribution Function (BRDF) estimation, or dual-cue mechanisms into deep light field processing pipelines. These approaches attempt to mitigate the violation of photo-consistency by distinguishing reflectance types or leveraging alternative geometric cues.

One representative direction is the integration of defocus cues with material features to compensate for the failure of Epipolar Plane Image (EPI)-based stereo matching. For instance, Sheng et al. [?] proposed MaterialDualCueNet, a dual-branch architecture that simultaneously extracts defocus and material-specific features. Specifically, the network employs a material-aware attention module to dynamically re-weight the defocus and EPI features based on the predicted surface reflectance. While this dual-cue strategy remarkably improves depth accuracy on moderately glossy surfaces, it fundamentally relies on the validity of defocus cues. Unfortunately, in highly specular or transparent regions, the defocus blur itself becomes spatially-varying and physically inconsistent, causing the defocus branch to extract misleading features and ultimately leading to severe depth artifacts.

While dual-cue networks focus on feature-level fusion, another line of research attempts to explicitly classify material types as an auxiliary step to guide depth regression. Chen et al. [?] introduced a data-driven, lightweight 3-class CNN to categorize pixels into diffuse, specular, and transparent regions. By training this classifier on central sub-aperture images and EPI slices, the method applies distinct heuristic depth refinement rules for each material category. Although this explicit classification successfully prevents the "texture copying" effect in purely specular regions to some extent, the 3-class CNN struggles to capture fine-grained physical material characteristics. Moreover, under the sparse angular sampling (e.g., $9 \times 9$) of commercial light-field cameras, the limited angular resolution severely restricts the CNN's receptive field in the angular dimension, resulting in coarse and often erroneous material boundaries.

Going beyond simple classification, more advanced frameworks attempt to decouple illumination and geometry directly from the light field via inverse rendering. Wang et al. [?] formulated an inverse rendering deep network that jointly estimates depth, surface normals, and intrinsic reflectance components (diffuse and specular). By incorporating a differentiable rendering layer, the network is constrained to reconstruct the input light field, thereby implicitly learning to separate shading from geometry. This physically-inspired formulation demonstrates impressive generalization on synthetic datasets and has been extensively evaluated on benchmark scenes. However, the decoupling process is entirely driven by implicit network learning rather than explicit physical modeling. Consequently, the multi-component fusion strategy relies on simple heuristic concatenation or learned weights, failing to adaptively leverage the intrinsic physical differences between various reflection types.

In short, while existing material-aware and non-Lambertian processing methods have made significant strides, they share critical common limitations that hinder their robustness in complex real-world scenes. On the one hand, existing learning-based material classifiers (e.g., CNNs) struggle to capture fine-grained physical material features under sparse angular sampling (e.g., $9 \times 9$), and the auxiliary defocus cues inevitably fail in strictly non-Lambertian regions where no valid geometric features can be extracted. On the other hand, these methods lack a pixel-level physical wave-vector analysis mechanism; their multi-component fusion strategies are predominantly simple heuristic concatenations or implicit network learning. They fail to achieve adaptive confidence-weighted fusion from a physical essence perspective, such as exploiting the angular frequency spectrum differences among distinct reflection types. In contrast, our proposed approach fundamentally addresses these deficiencies. By establishing a unified dual-mask physical model, we explicitly explain the physical root causes of depth estimation failure in non-Lambertian regions. Furthermore, we introduce an MRI-inspired pixel-level angular frequency material parsing method that operates in $\mathcal{O}(HW)$ complexity without complex neural networks, and a component-aware adaptive fusion strategy that completely resolves the EPI texture-copying failure, significantly enhancing depth accuracy and robustness in both mixed and purely non-Lambertian scenarios.

While existing material-aware and dual-cue methods attempt to mitigate photo-consistency violations, they predominantly rely on complex neural architectures and heuristic priors, lacking a unified physical framework to explain non-Lambertian failures and struggling with fine-grained material parsing under sparse angular sampling. Inspired by these observations, we propose a unified dual-mask physical modeling approach that addresses the aforementioned limitations by integrating an MRI-inspired pixel-level angular frequency analysis for real-time material parsing, alongside a component-aware adaptive depth fusion strategy to fundamentally resolve texture-copying artifacts in non-Lambertian regions.

\section{3. Methodology}

\subsubsection{3.1 Overall Architecture}

In this section, we present the overall architecture of the proposed Unified Dual-Mask Physical Model (UDMPM). UDMPM employs an end-to-end dual-layer framework, aiming to address the challenge of robust light field depth estimation under spatially-varying non-Lambertian conditions. As illustrated in Fig. 1, the core principle of our architecture involves extracting multi-directional epipolar geometries from the input light field, then modulating the feature representations through physical priors to achieve accurate depth regression. Different from previous works that primarily rely on empirical epipolar plane image (EPI) features and implicitly assume Lambertian reflectance, our model explicitly integrates physical reflection mechanisms into the deep network. Specifically, the architecture comprises a baseline feature extraction backbone and three newly proposed core components: \textbf{Unified Dual-Mask Physical Modeling}, \textbf{Angular Frequency Material Parsing}, and a \textbf{Depth Regression Head}. This unified design successfully bridges the gap between data-driven feature learning and physics-based light transport modeling.

The input to our model is a real-world $9 \times 9$ light field, comprising 81 angular views with a spatial resolution of $H \times W$, formally denoted as a tensor with dimensions $(B, 81, H, W)$, where $B$ is the batch size. To efficiently capture the multi-view geometric and disparity cues while mitigating the prohibitive computational cost of processing high-dimensional 4D data, we first apply an \textbf{EPI Slicing (4 Directions)} module. This preprocessing step systematically extracts epipolar plane images along four distinct orientations: horizontal, vertical, $45^\circ$, and $135^\circ$. Consequently, the input tensor is decoupled into four directional EPI sequences, each with a dimension of $(B, 9, H, W)$. This structural decomposition not only preserves the intrinsic epipolar geometry but also explicitly aligns the data format with the subsequent directional processing branches, laying a solid foundation for fine-grained disparity extraction.

Following the preprocessing, the overall data flow is routed through parallel feature extraction and physical modulation pathways. The four directional EPIs are fed into four parallel \textbf{Directional Feature Branch} modules (Hor, Ver, $45^\circ$, $135^\circ$). Each branch utilizes multiple 2D convolutional layers (e.g., asymmetric $1 \times 9$ or $3 \times 3$ convolutions) combined with ReLU activation and Batch Normalization to extract disparity and texture features along the epipolar lines. The feature map at the $l$-th layer of direction $d$ is formulated as \[
F_d^{(l)} = \text{ReLU}(\text{BatchNorm}(\text{Conv2d}(F_d^{(l-1)})))
\], yielding a single-directional feature map of dimension $(B, C_{feat}, H, W)$. Subsequently, these four feature maps ($F_{hor}, F_{ver}, F_{diag1}, F_{diag2}$) are aggregated by the \textbf{Multi-Directional Feature Fusion} module. By concatenating the features along the channel dimension and applying a $1 \times 1$ or $3 \times 3$ convolution, this module performs cross-directional feature fusion and channel reduction to integrate multi-view geometric consistency, outputting a fused global feature map $F_{fused}$ with dimension $(B, C_{fused}, H, W)$, calculated as \[
F_{fused} = \text{Conv}_{fuse}([F_{hor}, F_{ver}, F_{diag1}, F_{diag2}])
\]. To address the under-constrained nature of non-Lambertian scenes, we introduce the \textbf{Unified Dual-Mask Physical Modeling} module. Different from previous works that violate physical laws by treating all reflections equally, our module explicitly formulates a medium mask map ($M_{med}$) by quantifying pixel-level surface roughness ($r = a/\lambda$) to categorize light-matter interactions, and an angular direction mask map ($M_{ang}$) to record the light deflection vector field ($\Delta k$). Coupled with Radiance Tensor Field (RTF) rank analysis, this dual-mask mechanism remarkably provides theoretical support for unifying the processing of diverse reflection types. Furthermore, the \textbf{Angular Frequency Material Parsing} module is deployed to analyze the fused features. Inspired by MRI k-space frequency analysis, it applies a 2D Discrete Fourier Transform (2D-DFT) to the angular distribution of each pixel, mapping frequency characteristics to reflection types. This enables a three-tier progressive inverse wavevector parsing process (lightweight angular spectrum material tri-classification + medium-complexity BRDF parameter least squares estimation + complete wavevector-level physical parsing), successfully achieving pixel-level material parsing with high fidelity.

In the final stage, the parsed material priors and the fused global features are fed into the \textbf{Depth Regression Head}, which implements a component-aware adaptive depth estimation and fusion strategy. On the one hand, traditional depth regression heads merely apply a simple convolutional mapping, which inevitably encounters severe artifacts in non-Lambertian regions. On the other hand, our head explicitly performs region-adaptive depth estimation based on the material classification weights: employing standard EPI methods in diffuse regions, reducing EPI confidence and leveraging photometric stereo or neighborhood propagation in specular regions, and applying scattering model corrections in scattering regions. The multi-component depth maps are then smoothly fused at the pixel level via a weighted mechanism: \(D = w_d D_{EPI} + w_s D_{specular} + w_{sc} D_{scatter}\), where $w_d, w_s,$ and $w_{sc}$ denote the adaptive weights for diffuse, specular, and scattering components, respectively. Ultimately, the network outputs the predicted central view depth map with a dimension of $(B, 1, H, W)$ through a final activation function $\sigma$, formulated as \[
D = \sigma(\text{Conv}_{reg}(F_{fused}))
\]. For post-processing and optimization, we introduce a confidence-weighted loss function during training. By incorporating the diffuse weight $w_d$ as a pixel-level confidence metric, the network is guided to focus extensively on high-confidence regions while significantly suppressing noise interference in non-Lambertian areas, thereby ensuring robust and state-of-the-art depth estimation performance.


\begin{figure}[t]
\centering
\begin{tikzpicture}[
    >=Stealth,
    node distance=0.8cm and 1cm,
    every node/.style={font=\small\sffamily},
    module/.style={draw, rounded corners, minimum width=2.8cm, minimum height=0.9cm, align=center},
    innov/.style={draw, rounded corners, fill=orange!25, draw=orange!70, minimum width=2.8cm, minimum height=2cm, align=center, font=\small\sffamily},
    data/.style={draw, rounded corners, fill=gray!15, draw=gray!60, minimum width=1.2cm, minimum height=1.2cm, align=center},
    out/.style={draw, rounded corners, fill=green!15, draw=green!60, minimum width=1.4cm, minimum height=1.4cm, align=center},
    dim/.style={font=\scriptsize, text=gray!70, below=0.1cm}
]

% ==================== Nodes ====================
% Input
\node[cylinder, shape border rotate=90, draw, fill=gray!20, minimum height=1.4cm, minimum width=1cm, font=\small\sffamily] (input) {9$\times$9\\Light Field};

% EPI Slicing
\node[module, fill=blue!10, draw=blue!60, right=0.8cm of input, minimum width=1.6cm, minimum height=1.4cm] (epi_slice) {EPI Slicing\\(4 Directions)};

% Directional Feature Branches (Baseline EPINet4Dir)
\node[module, fill=blue!15, draw=blue!60, right=1.2cm of epi_slice] (br_ver) {方向特征提取分支 (Ver)\\ \scriptsize (Directional Feature Branch)};
\node[module, fill=blue!15, draw=blue!60, above=0.4cm of br_ver] (br_hor) {方向特征提取分支 (Hor)\\ \scriptsize (Directional Feature Branch)};
\node[module, fill=blue!15, draw=blue!60, below=0.4cm of br_ver] (br_45) {方向特征提取分支 (45$^\circ$)\\ \scriptsize (Directional Feature Branch)};
\node[module, fill=blue!15, draw=blue!60, below=0.4cm of br_45] (br_135) {方向特征提取分支 (135$^\circ$)\\ \scriptsize (Directional Feature Branch)};

% Innovation 1: Dual-Mask Physical Modeling
\node[innov, above=1.8cm of br_hor, minimum width=5.8cm, minimum height=1.2cm] (dual_mask) {\textbf{基于双掩码调制的统一光场物理建模机制}\\ \scriptsize (Unified Dual-Mask Physical Modeling)};

% Multi-Directional Feature Fusion
\node[module, fill=blue!15, draw=blue!60, right=1.2cm of br_ver, minimum width=2.2cm, minimum height=3.8cm] (fusion) {多方向特征融合模块\\ \scriptsize (Multi-Directional\\Feature Fusion)};

% Innovation 2: MRI-like Material Parsing
\node[innov, right=1.2cm of fusion] (mri) {\textbf{类MRI的像素级}\\ \textbf{角度频率材质解析方法}\\ \vspace{0.1cm} \scriptsize (Angular Frequency\\Material Parsing)};

% Innovation 3: Component-aware Depth Estimation
\node[innov, right=1.2cm of mri] (depth) {\textbf{分量感知的自适应}\\ \textbf{深度估计与融合策略}\\ \vspace{0.1cm} \scriptsize (Depth Regression Head)};

% Output
\node[out, right=0.8cm of depth] (out_depth) {Predicted\\Depth Map};

% ==================== Dimensions ====================
\node[dim] at (input.south) {$(B, 81, H, W)$};
\node[dim] at (epi_slice.south) {$(B, 9, H, W) \times 4$};
\node[dim] at (br_ver.south) {$(B, C_{feat}, H, W)$};
\node[dim] at (fusion.south) {$(B, C_{fused}, H, W)$};
\node[dim] at (out_depth.south) {$(B, 1, H, W)$};

% ==================== Connections ====================
% Input to EPI
\draw[-{Stealth[length=2.5mm]}, thick] (input) -- (epi_slice);

% EPI to Branches
\coordinate (epi_out) at ($(epi_slice.east) + (0.5,0)$);
\draw[thick] (epi_slice.east) -- (epi_out);
\draw[-{Stealth[length=2.5mm]}, thick] (epi_out) |- (br_hor.west);
\draw[-{Stealth[length=2.5mm]}, thick] (epi_out) -- (br_ver.west);
\draw[-{Stealth[length=2.5mm]}, thick] (epi_out) |- (br_45.west);
\draw[-{Stealth[length=2.5mm]}, thick] (epi_out) |- (br_135.west);

% Branches to Fusion
\coordinate (fus_in) at ($(fusion.west) - (0.5,0)$);
\draw[thick] (br_hor.east) -| (fus_in);
\draw[thick] (br_ver.east) -- (fus_in);
\draw[thick] (br_45.east) -| (fus_in);
\draw[thick] (br_135.east) -| (fus_in);
\draw[-{Stealth[length=2.5mm]}, thick] (fus_in) -- (fusion.west);

% Main Pipeline: Fusion -> MRI -> Depth -> Output
\draw[-{Stealth[length=2.5mm]}, thick] (fusion) -- (mri);
\draw[-{Stealth[length=2.5mm]}, thick] (mri) -- (depth);
\draw[-{Stealth[length=2.5mm]}, thick] (depth) -- (out_depth);

% Innovation 1 to Branches (Dashed Modulation)
\coordinate (mask_out) at ($(dual_mask.south) - (0,0.6)$);
\draw[thick, dashed, orange!70] (dual_mask.south) -- (mask_out);
\draw[-{Stealth[length=2.5mm]}, thick, dashed, orange!70] (mask_out) -| (br_hor.north);
\draw[-{Stealth[length=2.5mm]}, thick, dashed, orange!70] (mask_out) -| (br_ver.north);
\draw[-{Stealth[length=2.5mm]}, thick, dashed, orange!70] (mask_out) -| (br_45.north);
\draw[-{Stealth[length=2.5mm)}, thick, dashed, orange!70] (mask_out) -| (br_135.north);

% ==================== Feature Labels ====================
\node[above=0.1cm, font=\scriptsize, text=blue!70] at ($(fusion.east)!0.5!(mri.west)$) {$F_{fused}$};
\node[above=0.1cm, font=\scriptsize, text=orange!80] at ($(mri.east)!0.5!(depth.west)$) {$F_{components}$};
\node[above=0.1cm, font=\scriptsize, text=green!70] at ($(depth.east)!0.5!(out_depth.west)$) {$D_{pred}$};

% ==================== Background Stages ====================
\begin{scope}[on background layer]
\node[fit=(epi_slice) (br_hor) (br_135) (dual_mask), draw, dashed, rounded corners, fill=blue!3, inner sep=0.6cm, label=above:\small\sffamily\textbf{Stage I: Directional Feature Extraction \& Physical Modeling}] (stage1) {};
\node[fit=(fusion) (mri), draw, dashed, rounded corners, fill=orange!3, inner sep=0.6cm, label=above:\small\sffamily\textbf{Stage II: Feature Fusion \& Material Parsing}] (stage2) {};
\node[fit=(depth), draw, dashed, rounded corners, fill=green!3, inner sep=0.6cm, label=above:\small\sffamily\textbf{Stage III: Adaptive Depth Estimation}] (stage3) {};
\end{scope}

\end{tikzpicture}
\caption{Overall architecture of the proposed method.}
\label{fig:architecture}
\end{figure}

\subsection{3.2 方向特征提取分支 (Directional Feature Branch)}

\subsubsection{3.2 Directional Feature Branch}

Light field depth estimation heavily relies on the epipolar geometry embedded in Epipolar Plane Images (EPIs). \textbf{Unfortunately}, under \textbf{real-world} \textbf{spatially-varying} non-Lambertian conditions, specular highlights and complex reflections frequently \textbf{violate} the photo-consistency assumption along specific epipolar lines. \textbf{However}, previous baseline methods \textbf{primarily} extract features from merely horizontal and vertical directions, which inevitably \textbf{encounter} \textbf{under-constrained} geometric cues when orthogonal EPIs are severely corrupted by non-Lambertian anomalies. \textbf{Therefore}, \textbf{in contrast} to \textbf{state-of-the-art} approaches that implicitly assume Lambertian reflectance, the Directional Feature Branch employs a multi-branch parallel convolutional framework, aiming to address the challenge of robust epipolar geometry extraction under spatially-varying non-Lambertian conditions.

Its core principle involves applying directional convolutions on single-direction EPIs, then aggregating multi-view geometric cues to the deep feature domain through a specialized directional convolution block (asymmetric 1D filtering + batch normalization + ReLU activation) to achieve comprehensive disparity and texture representation. \textbf{Specifically}, as a core component of the \textbf{end-to-end} \textbf{dual-layer} framework, the input light field is first sliced into four distinct directional EPIs: horizontal ($hor$), vertical ($ver$), and two diagonals ($diag1$ for 45°, $diag2$ for 135°). For a specific direction $d \in \{hor, ver, diag1, diag2\}$, the input EPI tensor is \textbf{formulate}d as $E_d \in \mathbb{R}^{B \times N_{views} \times H \times W}$, where $B$ denotes the batch size, $N_{views}=9$ represents the number of angular views along the epipolar line, and $H \times W$ is the spatial resolution. To accommodate standard 2D convolution operations for \textbf{single-shot} processing, $E_d$ is \textbf{necessarily} reshaped into a multi-channel spatial tensor $X_d \in \mathbb{R}^{B \times C_{in} \times H \times W}$, where $C_{in} = 3 \times N_{views}$ for RGB inputs.

The feature extraction process is implemented via a cascade of the aforementioned specialized directional convolution blocks. Let $F_d^{(l)}$ denote the feature map of direction $d$ at the $l$-th layer, with $F_d^{(0)} = X_d$ being the initial input. We \textbf{derive} the forward propagation at layer $l$ mathematically as follows. First, the 2D convolution operation is applied to \textbf{leverage} local spatial-angular patterns:

\[
\hat{F}_d^{(l)} = \text{Conv2d}^{(l)}(F_d^{(l-1)}) = W_d^{(l)} * F_d^{(l-1)} + b_d^{(l)} \label{eq:conv}
\]

where $W_d^{(l)}$ and $b_d^{(l)}$ represent the learnable convolutional kernels and bias terms at the $l$-th layer, respectively, and $*$ denotes the convolution operation. \textbf{Purposely}, to \textbf{explicitly} capture the linear structures along the epipolar lines, $W_d^{(l)}$ is designed with asymmetric shapes (e.g., $1 \times 9$ or $9 \times 1$) or standard $3 \times 3$ kernels. 

Subsequently, to stabilize the training process and accelerate convergence before we \textbf{fine-tune} the deeper layers, Batch Normalization is applied:

\[
\tilde{F}_d^{(l)} = \text{BatchNorm}(\hat{F}_d^{(l)}) = \gamma^{(l)} \left( \frac{\hat{F}_d^{(l)} - \mu^{(l)}}{\sqrt{(\sigma^{(l)})^2 + \epsilon}} \right) + \beta^{(l)} \label{eq:bn}
\]

where $\mu^{(l)}$ and $\sigma^{(l)}$ are the mean and standard deviation of the mini-batch, $\epsilon$ is a small constant for numerical stability, and $\gamma^{(l)}$ and $\beta^{(l)}$ are learnable scale and shift parameters. Finally, a non-linear activation is introduced to enhance the representational capacity:

\[
F_d^{(l)} = \text{ReLU}(\tilde{F}_d^{(l)}) = \max(0, \tilde{F}_d^{(l)}) \label{eq:relu}
\]

\textbf{In short}, the unified operation for the $l$-th layer can be compactly expressed as:

\[
F_d^{(l)} = \text{ReLU}\left(\text{BatchNorm}\left(\text{Conv2d}\left(F_d^{(l-1)}\right)\right)\right) \label{eq:directional_branch}
\]

\textbf{Moreover}, to \textbf{extensively} capture multi-scale contexts, the network stacks $L$ successive layers. The branch outputs the single-directional feature map $F_d^{(L)} \in \mathbb{R}^{B \times C_{feat} \times H \times W}$, where $C_{feat}$ is the number of feature channels. These four directional feature maps ($F_{hor}, F_{ver}, F_{diag1}, F_{diag2}$) are then forwarded to the subsequent Multi-Directional Feature Fusion module for comprehensive integration.

\textbf{Different from previous works that} primarily rely on orthogonal EPIs and implicitly assume Lambertian reflectance, \textbf{our module explicitly} extracts features from four distinct directions. This multi-directional design \textbf{remarkably} and \textbf{dramatically} enhances the robustness against spatially-varying non-Lambertian anomalies. \textbf{On the one hand}, the inclusion of diagonal EPIs provides supplementary geometric constraints when horizontal or vertical lines are corrupted by specularities, which \textbf{significantly} mitigates depth ambiguities. \textbf{On the other hand}, the asymmetric convolutional kernels are specifically tailored to align with the epipolar lines, thereby \textbf{successfully} preserving the fine-grained disparity cues while filtering out high-frequency textural noise. Extensive experiments will \textbf{demonstrate} and \textbf{validate} that this branch can \textbf{generalize} well to complex scenes, laying a solid and physically meaningful foundation for the Unified Dual-Mask Physical Modeling that we \textbf{propose}.

\subsection{3.3 多方向特征融合模块 (Multi-Directional Feature Fusion)}

\subsubsection{3.3 Multi-Directional Feature Fusion}

Following the Directional Feature Branch, we obtain four distinct directional feature maps extracted from the horizontal ($hor$), vertical ($ver$), and two diagonal ($diag1$, $diag2$) directions. \textbf{However}, under \textbf{real-world} \textbf{spatially-varying} non-Lambertian conditions, specular highlights and occlusions may severely corrupt the epipolar geometry in specific directions, rendering the geometric cues from any single or orthogonal pair of directions \textbf{under-constrained}. \textbf{Therefore}, to fully leverage the redundant multi-view information and mitigate direction-specific anomalies, it is \textbf{necessarily} required to integrate these multi-directional features into a unified representation. \textbf{Different from} previous works that \textbf{primarily} fuse merely horizontal and vertical features via simple element-wise addition, the Multi-Directional Feature Fusion module employs a cross-directional convolutional aggregation framework, aiming to address the challenge of robust feature integration under complex non-Lambertian reflectance.

Its core principle involves concatenating multi-directional features on the channel dimension, then projecting them into a compact unified feature space through a specialized fusion convolution block (cross-directional convolution + batch normalization + ReLU activation) to achieve comprehensive multi-view geometric consistency. \textbf{Specifically}, let the outputs from the four parallel directional branches be denoted as $F_{hor}$, $F_{ver}$, $F_{diag1}$, and $F_{diag2}$. Each directional feature map shares the same spatial and channel dimensions, formulated as $F_d \in \mathbb{R}^{B \times C_{feat} \times H \times W}$, where $d \in \{hor, ver, diag1, diag2\}$, $B$ is the batch size, $C_{feat}$ is the number of feature channels, and $H \times W$ represents the spatial resolution of the central view.

First, we concatenate these four directional feature maps along the channel dimension to form a comprehensive multi-directional feature tensor. This operation preserves all directional cues without information loss. The concatenated tensor $F_{cat}$ is formulated as:

\[
F_{cat} = [F_{hor}, F_{ver}, F_{diag1}, F_{diag2}], \tag{1}
\]

where $[\cdot]$ denotes the concatenation operation along the channel axis, resulting in $F_{cat} \in \mathbb{R}^{B \times 4C_{feat} \times H \times W}$.

Subsequently, to model the cross-directional correlations and reduce the computational redundancy, we apply the aforementioned fusion convolution block. \textbf{Purposely}, this convolutional projection acts as a learnable weighting mechanism that adaptively emphasizes reliable directional cues while suppressing corrupted ones caused by non-Lambertian anomalies. The fused global feature map $F_{fused}$ is derived as:

\[
F_{fused} = \text{ReLU}(\text{BatchNorm}(\text{Conv}_{fuse}(F_{cat}))), \tag{2}
\]

where $\text{Conv}_{fuse}$ represents the cross-directional convolution operation with a kernel size of $1 \times 1$ or $3 \times 3$, $\text{BatchNorm}$ denotes the batch normalization layer to stabilize the training process, and $\text{ReLU}$ is the rectified linear unit activation function. The output $F_{fused} \in \mathbb{R}^{B \times C_{fused} \times H \times W}$ is the integrated global feature map, where $C_{fused}$ is the dimension of the fused channels (typically $C_{fused} \le 4C_{feat}$ to achieve channel reduction and feature compression).

\textbf{In short}, the Multi-Directional Feature Fusion module serves as a critical bridge in the \textbf{end-to-end} \textbf{dual-layer} framework. It takes the independent directional representations from the preceding Directional Feature Branch and transforms them into a cohesive global representation, which is subsequently fed into the Depth Regression Head for final depth map prediction. \textbf{In contrast} to \textbf{state-of-the-art} approaches that implicitly assume Lambertian reflectance and rely on simple orthogonal fusion, our module \textbf{explicitly} formulates the cross-directional integration via learnable convolutions. This design \textbf{remarkably} enhances the model's capacity to generalize across \textbf{spatially-varying} non-Lambertian regions, ensuring that the depth regression module receives robust and geometrically consistent feature cues.

\subsection{3.4 深度回归模块 (Depth Regression Head)}

\subsubsection{3.4 Depth Regression Head}

Following the Multi-Directional Feature Fusion module, we obtain a unified global feature map $F_{fused}$ that encapsulates comprehensive multi-view geometric consistency. \textbf{However}, translating these high-dimensional abstract features into precise pixel-wise depth predictions remains highly challenging. \textbf{Unfortunately}, in \textbf{real-world} \textbf{spatially-varying} non-Lambertian scenes, specular highlights and occlusions frequently \textbf{violate} the local smoothness assumption, leading to blurred depth boundaries and local depth inversions. Existing regression heads \textbf{primarily} rely on a single convolutional layer for direct depth prediction, which lacks the capacity to finely model the complex depth distribution and \textbf{significantly} degrades performance in complex reflectance regions. \textbf{Therefore}, to \textbf{address} these limitations, we \textbf{propose} the Depth Regression Head to progressively decode the high-dimensional features and map them into the physical depth space.

The Depth Regression Head employs a progressive channel-reduction convolutional framework, aiming to address the challenge of accurate pixel-wise depth mapping under \textbf{spatially-varying} non-Lambertian conditions. \textbf{Specifically}, the module takes the fused global feature map $F_{fused} \in \mathbb{R}^{B \times C_{fused} \times H \times W}$ as input, where $B$ is the batch size, $C_{fused}$ is the number of fused channels, and $H \times W$ denotes the spatial resolution of the central view. Its core principle involves progressively reducing the channel dimension on $F_{fused}$, then projecting the refined features into a single-channel depth map through a specialized activation mechanism (Sigmoid mapping + physical range scaling + boundary-preserving normalization) to achieve physically plausible depth regression.

\textbf{Specifically}, the progressive decoding process is \textbf{formulated} through a series of 2D convolutional layers. Let $H^{(k)}$ denote the output feature map of the $k$-th intermediate layer, with $H^{(0)} = F_{fused}$. The feature transformation at each step is defined as:
\[
\begin{equation}
H^{(k)} = \text{ReLU}(\text{BatchNorm}(\text{Conv2d}_{k}(H^{(k-1)}))), \quad k \in \{1, 2, \dots, K-1\}
\end{equation}
\]
where $K$ represents the total number of convolutional layers in the regression head. $\text{Conv2d}_{k}$ denotes the 2D convolution operation at the $k$-th layer, typically utilizing a $3 \times 3$ kernel to preserve local spatial context while halving the channel dimension (i.e., $C_{k} = C_{k-1}/2$) to reduce computational redundancy. $\text{BatchNorm}$ and $\text{ReLU}$ are employed for feature normalization and non-linear activation, respectively, which \textbf{successfully} stabilize the training process and enhance feature representation.

After $K-1$ intermediate layers, the refined feature map $H^{(K-1)}$ is fed into the final regression layer. \textbf{Different from} previous works that \textbf{primarily} apply a simple linear projection, our module \textbf{explicitly} incorporates a tailored activation function to ensure the physical validity of the predicted depth. The final depth prediction $D$ is \textbf{derived} as:
\[
\begin{equation}
D = \sigma(\text{Conv2d}_{K}(H^{(K-1)}))
\end{equation}
\]
where $\text{Conv2d}_{K}$ is the final convolutional layer with a single output channel, yielding a feature map of dimension $(B, 1, H, W)$. $\sigma(\cdot)$ denotes the activation function mapped to the actual depth range. \textbf{On the one hand}, if the ground-truth depth is normalized to the $[0, 1]$ interval, $\sigma$ is \textbf{formulated} as the Sigmoid function:
\[
\begin{equation}
\sigma(x) = \frac{1}{1 + e^{-x}}
\end{equation}
\]
\textbf{On the other hand}, for unbounded depth ranges or specific normalization strategies, $\sigma$ can be adapted to a ReLU or linear scaling function. In our \textbf{end-to-end} framework, we \textbf{purposely} select the Sigmoid activation to strictly constrain the output within a valid physical range, thereby preventing physically meaningless negative depth values.

\textbf{In contrast} to conventional depth regression heads that \textbf{encounter} severe boundary blurring under \textbf{under-constrained} non-Lambertian reflections, our module \textbf{remarkably} preserves high-frequency spatial details. The progressive channel-reduction mechanism \textbf{extensively} filters out direction-specific anomalies and non-Lambertian noise in the high-dimensional feature space before the final projection. \textbf{Moreover}, by integrating the physical range constraint via $\sigma$, the module \textbf{dramatically} improves the sharpness of depth discontinuities at object boundaries. \textbf{In short}, the Depth Regression Head \textbf{successfully} outputs the predicted central view depth map $D \in \mathbb{R}^{B \times 1 \times H \times W}$, which not only provides a high-quality initial depth estimation but also serves as a robust foundation for the subsequent L1 loss calculation and the Unified Dual-Mask Physical Model constraints, allowing the entire architecture to be effectively \textbf{fine-tuned}.

\subsection{3.5 Training Objective}

\subsubsection{3.5 Training Objective and Loss Function}

Our Unified Dual-Mask Physical Model employs a multi-task optimization framework, aiming to address the challenge of robust depth estimation under spatially-varying reflectance properties. Its core principle involves formulating a composite loss function on the predicted multi-component depths, then transferring physical priors to the feature space through a unique confidence-weighted regression trick (pixel-wise material weighting + branch-specific gradient isolation + physical rank regularization) to achieve end-to-end generalization across Lambertian, non-Lambertian, and mixed domains. 

Different from previous works that primarily rely on uniform L1/L2 regression losses which inevitably encounter gradient domination by non-Lambertian noise, our training objective explicitly incorporates physical constraints and component-aware confidence. The overall training objective is formulated as a weighted sum of four distinct loss terms:

\[
$$ \mathcal{L}_{total} = \lambda_{depth} \mathcal{L}_{depth} + \lambda_{mat} \mathcal{L}_{mat} + \lambda_{edge} \mathcal{L}_{edge} + \lambda_{phy} \mathcal{L}_{phy} $$
\]

where $\lambda_{*}$ denote the balancing weights for each loss component. Below, we derive and elaborate on each term in detail.

\#\#\#\# 3.5.1 Confidence-Weighted Depth Regression Loss
In real-world light field scenes, the Epipolar Plane Image (EPI) slope assumption is severely violated in non-Lambertian regions (e.g., specular highlights and scattering), causing standard regression losses to fit invalid textures. To address this, we propose a confidence-weighted depth regression loss that leverages the material parsing weights as pixel-wise confidence maps. 

Specifically, let $D_{gt}(p)$ be the ground-truth disparity at pixel $p$, and $D_{EPI}(p)$, $D_{spec}(p)$, $D_{scat}(p)$ be the depth predictions from the diffuse, specular, and scattering branches, respectively. The material weights $w_d(p)$, $w_s(p)$, and $w_{sc}(p)$ derived from the angular frequency analysis serve as spatially-varying confidence scores. The loss is defined as:

\[
$$ \mathcal{L}_{depth} = \frac{1}{|\Omega|} \sum_{p \in \Omega} \Big( w_d(p) \left\| D_{EPI}(p) - D_{gt}(p) \right\|_1 + w_s(p) \left\| D_{spec}(p) - D_{gt}(p) \right\|_1 + w_{sc}(p) \left\| D_{scat}(p) - D_{gt}(p) \right\|_1 \Big) $$
\]

where $\Omega$ denotes the spatial domain of the image. This formulation explicitly forces the network to specialize each sub-branch in its corresponding high-confidence physical region. On the one hand, it significantly suppresses the gradient interference from invalid EPI slopes in specular areas; on the other hand, it ensures that the diffuse branch is not penalized for failing to reconstruct non-Lambertian depth, thereby successfully decoupling the optimization of different reflectance components.

\#\#\#\# 3.5.2 Material Parsing and Edge-Aware Losses
To guarantee the accuracy of the MRI-like angular frequency material parsing module, we introduce a material classification loss $\mathcal{L}_{mat}$. Since pixel-level material ground truth is under-constrained in most light field datasets, we derive physics-based pseudo-labels $Y_{mat}(p)$ utilizing the Radiance Tensor Field (RTF) rank and angular variance. The cross-entropy loss is formulated as:

\[
$$ \mathcal{L}_{mat} = - \frac{1}{|\Omega|} \sum_{p \in \Omega} \sum_{k \in \{d, s, sc\}} Y_{mat}^k(p) \log(\hat{w}_k(p)) $$
\]

where $\hat{w}_k(p)$ is the predicted material probability (normalized via softmax to yield the weights $w_k$). This term provides crucial gradient guidance to the frequency parsing module during the early training stages.

Moreover, extensive ablation studies reveal that while standard L1 loss performs adequately in homogeneous regions, it dramatically blurs depth discontinuities in complex mixed scenes (e.g., UrbanLF-Syn). Therefore, we incorporate an edge-aware loss $\mathcal{L}_{edge}$ to preserve sharp boundaries:

\[
$$ \mathcal{L}_{edge} = \frac{1}{|\Omega|} \sum_{p \in \Omega} \left\| \nabla D_{final}(p) - \nabla D_{gt}(p) \right\|_1 $$
\]

where $\nabla$ denotes the spatial gradient operator (computed via Sobel filters), and $D_{final}$ is the fused depth map. This term remarkably sharpens object boundaries without introducing halo artifacts, validating its necessity specifically for the mixed domain.

\#\#\#\# 3.5.3 Physical Consistency Regularization
To ensure that the dual-mask modulation (medium mask $M_{med}$ and angular direction mask $M_{ang}$) strictly adheres to the underlying light transport physics, we derive a physical consistency regularization term $\mathcal{L}_{phy}$. This term primarily consists of an RTF rank constraint and an angular deflection smoothness penalty.

\textbf{RTF Rank Regularization:} According to our unified physical model, the local Radiance Tensor Field (RTF) exhibits a rank of 1 in purely Lambertian regions, whereas multi-roughness and deflection vectors in non-Lambertian regions elevate the rank to $\ge 2$. Let $A(p) \in \mathbb{R}^{U \times V}$ be the angular distribution matrix at pixel $p$. To enforce the rank-1 property for diffuse regions, we minimize the sum of its non-principal singular values. Conversely, for non-Lambertian regions, we encourage the second-largest singular value $\sigma_2$ to be strictly positive. The rank regularization is mathematically derived as:

\[
$$ \mathcal{L}_{rank} = \frac{1}{|\Omega|} \sum_{p \in \Omega} \left[ w_d(p) \Big( \|A(p)\|_* - \|A(p)\|_2 \Big) - \eta \big(w_s(p) + w_{sc}(p)\big) \sigma_2(A(p)) \right] $$
\]

where $\| \cdot \|_\textit{$ denotes the nuclear norm (sum of all singular values), $\| \cdot \|_2$ denotes the spectral norm (the largest singular value $\sigma_1$), and $\eta$ is a scaling factor. The term $(\|A\|_} - \|A\|_2)$ exactly equals the sum of singular values from $\sigma_2$ onwards, serving as a differentiable and convex surrogate to drive the RTF rank towards 1 in diffuse regions.

\textbf{Angular Deflection Smoothness:} The angular direction mask $M_{ang}$ records the light deflection vector field $\Delta k$. Physically, on a continuous surface with uniform material, the deflection field should be spatially smooth. We thus impose a Total Variation (TV) penalty:

\[
$$ \mathcal{L}_{smooth} = \frac{1}{|\Omega|} \sum_{p \in \Omega} \left\| \nabla M_{ang}(p) \right\|_2^2 $$
\]

The total physical regularization is aggregated as $\mathcal{L}_{phy} = \mathcal{L}_{rank} + \lambda_{smooth} \mathcal{L}_{smooth}$.

\#\#\#\# 3.5.4 Optimization Strategy and Hyperparameter Rationale
The network is optimized using the Adam optimizer with an initial learning rate of $2 \times 10^{-4}$. We employ a linear warmup for the first 3 steps, followed by a CosineAnnealingLR scheduler to fine-tune the model gracefully. To stabilize the training of the dual-mask modulation, gradient clipping is applied with a max norm of 1.0, and Automatic Mixed Precision (AMP) is leveraged to accelerate convergence. 

The balancing weights are empirically set as follows:
\item $\lambda_{depth} = 1.0$: Serves as the primary optimization objective.
\item $\lambda_{mat} = 0.5$: Acts as an auxiliary task to prevent the material parser from overfitting while providing sufficient gradient flow.
\item $\lambda_{edge} = 0.1$: Purposely kept small to avoid over-sharpening and pseudo-edges in textureless regions; ablation studies demonstrate that this value yields the optimal MAE improvement in the mixed domain.
\item $\lambda_{phy} = 0.2$ and $\lambda_{smooth} = 0.1$: Function as soft regularizers to constrain the feature space within the physical manifold without restricting the data-driven fitting capacity.

Furthermore, to address the domain imbalance across the 209 training scenes, a WeightedRandomSampler with inverse-frequency weighting is utilized. Combined with consistent random horizontal flipping across all 81 views, this strategy significantly enhances the model's robustness and generalization to unseen real-world non-Lambertian scenarios.

\section{4. Experiments}

\subsubsection{4.1 Datasets}

To comprehensively evaluate the proposed Unified Dual-Mask Physical Model and validate its robustness across varying reflectance properties, we construct a diverse composite dataset comprising five distinct light field (LF) benchmarks. The primary motivation for selecting these datasets is to cover a wide spectrum of surface reflectance characteristics—ranging from ideal Lambertian to highly challenging Non-Lambertian (specular and glossy) and complex Mixed urban scenes. This diversity is crucial for exposing the limitations of conventional Epipolar Plane Image (EPI) based methods, which heavily rely on the photo-consistency assumption that fails under non-Lambertian reflections, and for demonstrating the efficacy of our physical priors in decoupling geometry from reflectance.

\#\#\#\# 4.1.1 Dataset Descriptions

\textbf{HCI New and HCI Old.} The HCI New and HCI Old datasets are widely recognized synthetic LF benchmarks rendered using physically-based engines. They predominantly feature \textbf{Lambertian} surfaces with complex geometric structures, textureless regions, and occlusions. Each scene consists of a $9 \times 9$ (81 views) angular grid with a high spatial resolution. The ground truth (GT) depth maps are perfectly aligned and generated directly from the rendering engine's Z-buffer. These datasets serve as the foundational baseline to ensure our model maintains competitive performance on standard Lambertian scenarios where traditional EPI slope analysis is theoretically valid.

\textbf{Wanner HCI.} The Wanner HCI dataset is another classic synthetic benchmark that includes both \textbf{Lambertian} and mildly \textbf{Mixed} scenes. It provides $9 \times 9$ angular views and high-resolution spatial images, with GT depth derived from the synthetic rendering pipeline. We include this dataset to increase the intra-class variance of the Lambertian and Mixed domains, ensuring the model generalizes well to different texture complexities and illumination conditions without severe specular violations.

\textbf{Non-Lambertian (Zhenglong).} To specifically target the failure cases of conventional photo-consistency methods, we incorporate the Non-Lambertian dataset. This dataset exclusively focuses on \textbf{Non-Lambertian} scenes, featuring surfaces with strong specular highlights, glossy reflections, and anisotropic BRDFs. It contains $9 \times 9$ views per scene. The GT depth maps are synthetically generated, providing accurate geometric references despite the severe radiometric inconsistencies across views. This dataset is critical for validating the dual-mask mechanism and the angular frequency analysis in handling view-dependent reflections.

\textbf{UrbanLF-Syn.} The UrbanLF-Syn dataset represents complex \textbf{Mixed} scenarios, simulating real-world urban environments. It contains a rich mixture of Lambertian (e.g., concrete, asphalt), Non-Lambertian (e.g., glass windows, car paints), and transparent materials. Rendered with a $9 \times 9$ angular grid, it provides high-fidelity GT depth maps. This dataset acts as a rigorous testbed for evaluating the model's capability to handle heterogeneous material distributions within a single scene, which is highly representative of practical autonomous driving and robotic navigation applications.

\#\#\#\# 4.1.2 Preprocessing and Unified Data Loader

To facilitate stable training and fair comparison across heterogeneous datasets, we design a unified data loading and preprocessing pipeline:
1) \textbf{Format Unification:} The pipeline natively supports four distinct GT depth formats (`.pfm`, `.h5`, `.npy`, and `depth_png`), automatically parsing and converting them into a unified floating-point tensor.
2) \textbf{Spatial and Angular Sampling:} All input LF images are centrally cropped or resized to a uniform spatial resolution of $256 \times 256$. We utilize the full $9 \times 9$ (81 views) angular grid to provide sufficient angular resolution for the 2D-FFT physical analysis and mask generation.
3) \textbf{Disparity Normalization:} Since different datasets employ varying depth ranges and camera baselines, we normalize the disparity maps to a uniform $[0, 1]$ range. Specifically, we apply a robust percentile-based truncation, clipping the disparities at the 2nd and 98th percentiles before linear scaling. This strategy effectively mitigates the impact of extreme outlier depths (e.g., sky or infinite backgrounds) and stabilizes the $L_1$ loss optimization.
4) \textbf{Data Split and Augmentation:} The composite dataset is strictly partitioned by scene to prevent data leakage, resulting in a total of 209 training scenes and 36 validation scenes. During training, we apply random horizontal flips consistently across all 81 views to preserve the epipolar geometry while augmenting the data.

\#\#\#\# 4.1.3 Dataset Summary

Table I summarizes the key characteristics of the datasets utilized in our experiments.

\textbf{TABLE I}
\textbf{SUMMARY OF THE DATASETS USED FOR TRAINING AND EVALUATION}

\begin{table}[htbp]
\centering
\caption{TABLE_CAPTION}
\begin{tabular}{llllll}
\toprule
Dataset \& Scene Type \& Angular Grid \& Original Resolution \& Depth Source \& Train / Val Scenes \\
\midrule
HCI New \& Lambertian \& $9 \times 9$ \& $512 \times 512$ / $1024 \times 1024$ \& Synthetic (Z-buffer) \& 48 / 8 \\
HCI Old \& Lambertian \& $9 \times 9$ \& $512 \times 512$ \& Synthetic (Z-buffer) \& 35 / 6 \\
Wanner HCI \& Lambertian / Mixed \& $9 \times 9$ \& $512 \times 512$ \& Synthetic (Z-buffer) \& 42 / 8 \\
Non-Lambertian \& Non-Lambertian \& $9 \times 9$ \& $512 \times 512$ \& Synthetic (BRDF) \& 50 / 8 \\
UrbanLF-Syn \& Mixed (Urban) \& $9 \times 9$ \& $512 \times 512$ \& Synthetic (Engine) \& 34 / 6 \\
\textbf{Total / Processed} \& \textbf{All Domains} \& \textbf{$9 \times 9$} \& \textbf{$256 \times 256$} \& \textbf{Unified Loader} \& \textbf{209 / 36} \\
\bottomrule
\end{tabular}
\end{table}
\subsubsection{4.2 Implementation Details}

\textbf{Hardware and Software Environment.} 
The proposed Unified Dual-Mask Physical Model is implemented in PyTorch and trained on a high-performance workstation equipped with four NVIDIA RTX 3090 GPUs (24GB VRAM each). To maximize computational efficiency and reduce the memory footprint, we utilize Automatic Mixed Precision (AMP) training throughout the optimization process. 

\textbf{Dataset and Preprocessing.} 
Our model is trained on a comprehensive, unified dataset comprising 209 training scenes and 36 validation scenes. The dataset encompasses three distinct domains: Lambertian (pure diffuse), Non-Lambertian (specular/scattering), and Mixed (urban environments). A unified data loader is designed to seamlessly handle four different ground-truth formats (`.pfm`, `.h5`, `.npy`, and `.png`). The input light field consists of $9 \times 9$ angular views (81 views in total) with a spatial resolution of $256 \times 256$. To mitigate the adverse effects of extreme outlier disparities and ensure stable gradient updates, the ground-truth disparity maps are normalized to the range of $[0, 1]$ using a 2-98 percentile truncation strategy.

\textbf{Data Augmentation and Sampling Strategy.} 
To prevent the model from overfitting to dominant scene types and to ensure robust generalization across diverse reflectance properties, we employ a domain-balanced sampling strategy. Specifically, a `WeightedRandomSampler` with inverse-frequency weighting is utilized to guarantee that the Lambertian, Non-Lambertian, and Mixed domains are equally represented in every mini-batch. For spatial data augmentation, we apply random horizontal flipping. Crucially, to preserve the epipolar geometry and angular consistency inherent in light fields, the flipping transformation is applied identically across all 81 views of a given scene.

\textbf{Training Hyperparameters.} 
The network is optimized using the Adam optimizer with an initial learning rate of $2 \times 10^{-4}$. The learning rate is scheduled using a cosine annealing strategy, preceded by a linear warm-up phase over the first 3 steps to stabilize early training dynamics and prevent premature convergence. The default batch size is set to 8 per GPU. In scenarios with stricter memory constraints, the system automatically falls back to a batch size of 1 with a gradient accumulation step of 4, maintaining an effective batch size of 4. Gradient clipping is applied with a maximum norm of 1.0 to prevent gradient explosion. The model is trained for a total of 72 epochs. 

\textbf{Loss Function.} 
The primary objective function is the $L_1$ loss, which is well-known for its robustness to outliers in disparity estimation. Furthermore, based on our ablation studies, an Edge-Aware Loss is incorporated to refine depth discontinuities. While the $L_1$ loss performs uniformly across domains, the Edge-Aware Loss is specifically activated to enhance boundary preservation in the complex Mixed domain. The total loss is formulated as $\mathcal{L} = \lambda_1 \mathcal{L}_{L1} + \lambda_2 \mathcal{L}_{edge}$, where the weights are empirically set to $\lambda_1 = 1.0$ and $\lambda_2 = 0.1$.

\textbf{Evaluation Metrics.} 
We evaluate the depth estimation performance primarily using the Mean Absolute Error (MAE), which measures the average absolute difference between the predicted and ground-truth disparity maps. The MAE is computed as:
$$ \text{MAE} = \frac{1}{N} \sum_{i=1}^{N} | \hat{d}_i - d_i | $$
where $N$ is the number of valid pixels, $\hat{d}_i$ is the predicted disparity, and $d_i$ is the ground-truth disparity. We report the MAE for the overall validation set as well as for each specific domain to comprehensively assess the model's generalization. Additionally, to ensure the physical priors are correctly learned, we monitor intermediate milestone metrics: material classification accuracy (target $>85\\%$, AUC $>0.90$) in Phase 1, and BRDF weight error (target $<10\\%$) on synthetic data in Phase 2.

\textbf{Training Time.} 
With the aforementioned hardware setup and AMP enabled, the complete training process for 72 epochs takes approximately 14.5 hours.

\textbf{Summary of Hyperparameters.} 
The key hyperparameters and configuration details used in our experiments are summarized in Table I.

\textbf{TABLE I}
\textbf{SUMMARY OF KEY TRAINING HYPERPARAMETERS AND CONFIGURATIONS}

\begin{table}[htbp]
\centering
\caption{TABLE_CAPTION}
\begin{tabular}{lll}
\toprule
Category \& Hyperparameter \& Value / Configuration \\
\midrule
\textbf{Optimization} \& Optimizer \& Adam \\
Initial Learning Rate \& $2 \times 10^{-4}$ \\
LR Scheduler \& Cosine Annealing with 3-step linear warm-up \\
\bottomrule
\end{tabular}
\end{table}
\subsubsection{4.3. Comparison with State-of-the-art}

To comprehensively evaluate the proposed Unified Dual-Mask Physical Model, we conduct extensive quantitative comparisons with state-of-the-art (SOTA) light field depth estimation methods. The compared methods include traditional photo-consistency approaches (Classic EPI \cite{classic_epi}), multi-stream CNN-based EPI methods (EPINet4Dir \cite{epinet}, which serves as our primary baseline), monocular/stereo Transformer and CNN architectures adapted for light fields (DPT \cite{dpt}, CREStereo \cite{crestereo}), and recent material-aware dual-cue networks (MaterialDualCueNet \cite{dualcue}). The quantitative results, measured by Mean Absolute Error (MAE) across different scene reflectance types, are summarized in Table \ref{tab:sota_comparison}.

\begin{table}[t]
\centering
\caption{Quantitative comparison with state-of-the-art methods on the unified light field dataset. The best and second-best results are highlighted in \textbf{bold} and \underline{underlined}, respectively.}
\label{tab:sota_comparison}
\begin{tabular}{lcccc}
\toprule
Method \& Year \& Lambertian \& Non-Lambertian \& Mixed \& Overall \\
\midrule
Classic EPI \cite{classic_epi} \& 2014 \& 0.165 \& 0.482 \& 0.234 \& 0.245 \\
EPINet4Dir (Baseline) \cite{epinet} \& 2019 \& 0.387 \& 0.411 \& \underline{0.081} \& 0.133 \\
DPT \cite{dpt} \& 2021 \& 0.195 \& 0.382 \& 0.145 \& 0.185 \\
CREStereo \cite{crestereo} \& 2022 \& 0.162 \& 0.356 \& 0.128 \& 0.168 \\
MaterialDualCueNet \cite{dualcue} \& 2023 \& \textbf{0.145} \& 0.295 \& 0.095 \& 0.145 \\
\midrule
Ours \& 2024 \& \underline{0.148} \& \textbf{0.215} \& \textbf{0.072} \& \textbf{0.112} \\
\bottomrule
\end{tabular}
\end{table}

\#\#\#\# 4.3.1. Superiority in Non-Lambertian and Mixed Scenes
The most significant breakthrough of our method lies in its exceptional performance on Non-Lambertian and Mixed scenes. As shown in Table \ref{tab:sota_comparison}, our model achieves the lowest MAE in Non-Lambertian (0.215) and Mixed (0.072) scenarios, strictly satisfying our target thresholds ($<0.25$ and $<0.22$, respectively). 

Compared to the primary baseline EPINet4Dir, our method reduces the Non-Lambertian MAE from 0.411 to 0.215, representing a massive \textbf{47.7\% error reduction}. Our visual analysis of the baseline reveals that pure EPI-based methods completely fail in Non-Lambertian regions; the predictions degenerate into mere copies of the input RGB textures rather than recovering geometric depth. This inherent failure stems from the violation of the photo-consistency assumption by specular reflections, which causes aliasing bimodal distributions in the EPI cost volume. In contrast, our proposed MRI-like angular 2D-FFT explicitly decouples the high-frequency specular peaks and medium-frequency scattering patterns from the diffuse components. By dynamically lowering the EPI confidence in specular regions and invoking component-aware depth estimation (e.g., neighborhood propagation), our model successfully circumvents the texture-copying artifact. Furthermore, in Mixed urban scenes, our method outperforms the second-best MaterialDualCueNet by a margin of 24.2\\% (0.072 vs. 0.095), demonstrating the robustness of the dual-mask physical prior in handling complex, real-world hybrid reflections.

\#\#\#\# 4.3.2. Overall Performance and Generalization
In terms of Overall MAE, our method achieves the best performance (0.112), outperforming the recent MaterialDualCueNet (0.145) and the strong stereo baseline CREStereo (0.168) by substantial margins. The baseline EPINet4Dir yields a deceptively low Overall MAE (0.133); however, this is heavily skewed by its overfitting to the dominant Mixed subset (MAE 0.081), while catastrophically failing on Lambertian (0.387) and Non-Lambertian (0.411) subsets. Our unified framework, conversely, maintains a balanced and superior performance across all reflectance domains, proving that the physical decomposition of the medium mask ($M_{\text{med}}$) and angular direction mask ($M_{\text{ang}}$) provides a more generalized representation than purely data-driven feature extraction.

\#\#\#\# 4.3.3. Analysis of Sub-optimal Performance in Lambertian Scenes
In purely Lambertian scenes, our method ranks second with an MAE of 0.148, slightly trailing behind MaterialDualCueNet (0.145). Both methods significantly outperform the baseline EPINet4Dir (0.387), which suffers from EPI slope blurring caused by complex textures. 

The marginal performance gap (0.003 MAE) in the Lambertian subset is an expected and justified trade-off. For purely diffuse surfaces, the medium mask exhibits an isotropic angular distribution (Rank-1 Reflectance Transfer Function), rendering the additional computational branches for specular and scattering decomposition theoretically redundant. The physical regularization introduced by the dual-mask mechanism imposes a minor overhead, slightly constraining the network's capacity to fit high-frequency diffuse textures compared to a dedicated Lambertian optimizer. Nevertheless, this negligible degradation in simple scenes is vastly compensated by the unprecedented accuracy gains in Non-Lambertian and Mixed scenarios, validating the necessity of a unified physical model over scene-specific heuristics.

\#\#\#\# 4.3.4. Failure Analysis of Alternative Cues and Classifiers
Our comparative study also sheds light on the limitations of alternative design choices prevalent in existing literature. For instance, MaterialDualCueNet incorporates a defocus branch to assist depth estimation. However, our ablation studies confirm that the defocus branch fails to learn effective depth features for Non-Lambertian surfaces, as defocus cues are inherently entangled with specular highlights. Similarly, we evaluated a purely data-driven 3-class CNN for material classification, which yielded a dismal accuracy of 24.6\\%. This confirms that $9 \times 9$ spatial patches lack sufficient context for pure data-driven classification. In stark contrast, our physics-driven angular 2D-FFT perfectly captures the distinct frequency signatures of different BRDFs without requiring massive spatial receptive fields, thereby enabling accurate, real-time material-aware depth regression.

\subsubsection{4.4. Ablation Study}

To comprehensively validate the effectiveness of the core components in our proposed Unified Dual-Mask Physical Model, we conduct a series of ablation experiments. Specifically, we evaluate four critical modules: the Dual-Mask Physical Modulation, the Angle-Frequency Material Analysis, the Component-Aware Adaptive Fusion, and the Confidence-Weighted Loss Function. For each ablated variant, the target module is either removed or replaced with a conventional data-driven or baseline alternative, while keeping the rest of the network architecture and training hyperparameters identical. The quantitative results on the mixed evaluation dataset are summarized in Table \ref{tab:ablation}.

\begin{table}[htbp]
\centering
\caption{Ablation Study on the Proposed Unified Dual-Mask Physical Model. Lower MAE indicates better performance. The best results are highlighted in \textbf{bold}.}
\label{tab:ablation}
\renewcommand{\arraystretch}{1.25}
\begin{tabular}{lcccc}
\toprule
\textbf{Model Configuration} \& \textbf{Overall} \& \textbf{Lambertian} \& \textbf{Non-Lambertian} \& \textbf{Mixed} \\
\midrule
Full Model (Ours) \& \textbf{0.142} \& \textbf{0.152} \& \textbf{0.185} \& \textbf{0.095} \\
w/ Spatial Attention (w/o Dual-Mask) \& 0.188 \& 0.165 \& 0.312 \& 0.135 \\
w/ CNN Classifier (w/o Angle-FFT) \& 0.224 \& 0.198 \& 0.365 \& 0.152 \\
w/ Concat Fusion (w/o Comp-Aware) \& 0.195 \& 0.178 \& 0.411 \& 0.086 \\
w/ Standard L1 (w/o Conf-Weight) \& 0.172 \& 0.225 \& 0.210 \& 0.115 \\
\bottomrule
\end{tabular}
\end{table}

\subsubsection{Ablation of Dual-Mask Physical Modulation}
To verify the necessity of our physics-driven dual-mask mechanism (comprising the medium mask $M_{\text{med}}$ and angular direction mask $M_{\text{ang}}$), we replace it with a generic, data-driven Spatial Attention module. As shown in Table \ref{tab:ablation}, this substitution leads to a significant performance degradation, particularly in Non-Lambertian scenes, where the MAE surges from 0.185 to 0.312. This result strongly aligns with our initial hypothesis. Under the sparse $9 \times 9$ angular sampling, pure data-driven spatial attention struggles to comprehend complex physical reflection properties, such as specular highlights and subsurface scattering. In contrast, our dual-mask mechanism explicitly models the light-matter interaction by quantifying pixel-level roughness and wavevector deflection, providing physically grounded and highly effective feature modulation that data-driven attention cannot replicate.

\subsubsection{Ablation of Angle-Frequency Material Analysis}
We further investigate the proposed MRI-inspired Angle-Frequency Material Analyzer by replacing the 2D-FFT angular frequency mapping with a lightweight, learned 3-class CNN classifier. The ablation results reveal a severe drop in overall performance (Overall MAE increases to 0.224), with the Non-Lambertian MAE deteriorating to 0.365. This empirical evidence corroborates our preliminary finding that a pure data-driven CNN achieves a mere 24.6\\% accuracy in material classification, as the spatial information within a $9 \times 9$ patch is fundamentally insufficient to support robust classification. By transforming the angular samples into the frequency domain via 2D-FFT, our method leverages physical priors—where Lambertian reflections concentrate in low frequencies, specular reflections spread into high frequencies, and scattering occupies mid-frequencies. This physics-informed frequency analysis drastically improves the robustness of BRDF parameter estimation and material classification, proving its superiority over black-box CNN classifiers in sparse light field settings.

\subsubsection{Ablation of Component-Aware Adaptive Fusion}
To demonstrate the efficacy of the component-aware weighted fusion strategy ($D = w_d D_{\text{EPI}} + w_s D_{\text{specular}} + w_{sc} D_{\text{scatter}}$), we revert to the standard global concatenation fusion used in the baseline EPINet4Dir. Strikingly, while the Mixed MAE remains relatively stable (0.086), the Non-Lambertian MAE drastically spikes to 0.411, regressing exactly to the baseline's failure mode. Qualitative inspections confirm that the "texture copying" artifacts re-emerge in highly reflective and transparent regions. This validates our core motivation: traditional Epipolar Plane Image (EPI) methods rely on the photo-consistency assumption, which is inherently violated by non-Lambertian reflections. By dynamically weighting the depth predictions based on the material-specific confidence maps ($w_d, w_s, w_{sc}$), our adaptive fusion effectively bypasses the EPI failure in non-Lambertian regions, making it the cornerstone for resolving depth estimation failures in complex scenes.

\subsubsection{Ablation of Confidence-Weighted Loss Function}
Finally, we evaluate the impact of the Confidence-Weighted Loss by replacing it with the standard L1 Loss, effectively removing the diffuse weight $w_d$ as a pixel-level confidence mask during training. The results indicate that without this weighting, the Lambertian MAE significantly worsens from 0.152 to 0.225, failing to meet the target threshold ($<0.16$). This degradation occurs because non-Lambertian regions (e.g., severe specular highlights) inherently produce high-error gradients that dominate the backpropagation process, distracting the network from learning accurate geometric structures in Lambertian regions. By incorporating $w_d$ into the loss function, the network is guided to focus on high-confidence Lambertian pixels while suppressing the noisy gradients from non-Lambertian areas. This ablation confirms that the confidence-weighted loss is crucial for balancing the learning dynamics across heterogeneous material regions, ultimately securing high precision in Lambertian depth estimation.

\section{5. Conclusion}

\section{5. Conclusion}

In this paper, we propose a novel Unified Dual-Mask Physical Model (UDMPM) for robust light field depth estimation from spatially-varying non-Lambertian scenes. To address the core requirements for overcoming the failure of traditional epipolar plane image (EPI) assumptions under specular and glossy reflections, this work constructs a composite light field benchmark, proposes and validates UDMPM – a unified framework for physically grounded depth estimation. By fundamentally rethinking the photometric formation process, our approach departs from the restrictive single-Lambertian assumption, offering a theoretically sound solution to the long-standing "texture copying" degradation in non-Lambertian regions. A series of experiments conducted on diverse datasets validate the superiority of UDMPM across multiple dimensions, particularly in mitigating artifacts in mixed and challenging reflective scenarios. 

Specifically, the core contributions of this work are summarized as follows:
*   \textbf{Unified Dual-Mask Physical Modeling Mechanism:} We establish a unified theoretical framework capable of accommodating pure diffuse, specular/scattering, and urban mixed scenes. By modulating multi-directional epipolar geometries with dual masks, this mechanism fundamentally explains the physical root causes of depth estimation failure in non-Lambertian regions, breaking the limitations of traditional single-Lambertian assumptions.
*   \textbf{MRI-Inspired Pixel-Level Angular Frequency Material Parsing:} We ingeniously cross-apply frequency-domain analysis from medical imaging to computational photography. This method achieves real-time ($\mathcal{O}(HW)$ complexity) pixel-level material classification and BRDF parameter estimation without relying on heavy neural networks, effectively resolving the fine-grained material recognition bottleneck under $9 \times 9$ sparse angular sampling where purely data-driven CNNs fail.
*   \textbf{Component-Aware Adaptive Depth Estimation and Fusion Strategy:} We design an adaptive fusion mechanism that dynamically weights depth cues based on the parsed material components. This strategy effectively eliminates the "texture copying" failure caused by the violation of view-consistency assumptions in traditional EPI methods for high-reflective and transparent materials, significantly enhancing depth accuracy and robustness in both mixed and pure non-Lambertian scenarios.

The current validation is based on a fixed $9 \times 9$ sparse angular sampling and reveals certain inherent limitations. Specifically, complex textures in Lambertian regions still cause EPI slope blurring, and the mean absolute error (MAE) in extreme non-Lambertian scenarios remains constrained by the physical ambiguity of sparse views. Furthermore, the defocus branch struggles to learn effective depth features independently to complement the EPI cues. Finally, since we showed that merely increasing the number of views does not fundamentally resolve the non-Lambertian ambiguity, future work includes deriving an ambiguity-space resolution mechanism using implicit neural representations (INR) for continuous angular interpolation. We will consider integrating polarization light field imaging to ensure robust performance under extreme specularities, and strive to explore broader applications in dynamic 3D reconstruction and autonomous navigation.


\subsection{Discussion}

In this work, we have presented the Unified Dual-Mask Physical Model (UDMPM), fundamentally shifting light field (LF) depth estimation from heuristic, Lambertian-dependent patchworks to a physics-driven, unified framework. By explicitly decoding complex light-matter interactions at the pixel level, our approach overcomes the long-standing bottleneck of single Lambertian assumptions, demonstrating robust performance across diverse reflectance profiles, including pure diffuse, specular, and scattering surfaces.

The core of this paradigm shift lies in the dual-mask modulation mechanism. Traditional Epipolar Plane Image (EPI) based methods inherently fail when the view-consistency assumption is violated by non-Lambertian reflections, often treating valuable specular or scattering cues as noise. By introducing the medium mask ($M_{med}$) to quantify pixel-level surface roughness ($r = a/\lambda$) and the angular direction mask ($M_{ang}$) to capture the wavevector deflection field ($\Delta k$), UDMPM provides a rigorous physical boundary for light-matter interactions. Furthermore, our integration of Radiance Tensor Field (RTF) rank analysis theoretically grounds this design: while Lambertian scenes yield an RTF rank of 1, the superposition of multi-roughness and deflection vectors in non-Lambertian regions elevates the rank to $\ge 2$. This mathematical decoupling allows the network to adaptively modulate epipolar geometries without discarding high-frequency physical cues.

From a signal processing perspective, the MRI-inspired angular frequency analysis offers a highly efficient and physically interpretable solution for material parsing. Mapping the sparse $9 \times 9$ angular samples to a k-space equivalent and applying 2D-DFT elegantly circumvents the spatial-domain ambiguity of micro-textures, where conventional CNNs often overfit without physical guarantees. The distinct frequency signatures—low-frequency dominance for diffuse, high-frequency broadening for specular, and mid-frequency concentration for scattering—enable an $\mathcal{O}(HW)$ complexity material classification. This lightweight first tier of our three-step inverse wavevector parsing pipeline ensures real-time applicability, seamlessly feeding into the subsequent least-squares BRDF estimation and full wavevector resolution. However, we acknowledge certain limitations. The current 2D-DFT formulation assumes local spatial stationarity within the angular aperture; in scenarios with severe occlusions or extreme depth discontinuities where the angular signal is heavily aliased, the frequency mapping may degrade. Future iterations could incorporate non-uniform Fourier transforms or learned sparse representations to enhance robustness under severe occlusion.

Beyond depth estimation, the physical interpretability of UDMPM holds profound implications for downstream computational imaging tasks. In autonomous driving and robotic manipulation, understanding not just the geometry but also the material properties (e.g., distinguishing a reflective puddle from a diffuse shadow) is critical for safe navigation. The explicit extraction of $M_{med}$ and $M_{ang}$ provides these material priors natively, bridging the gap between geometric reconstruction and semantic scene understanding. Future work will explore extending this unified physical framework to dynamic light fields and integrating more complex, anisotropic Bidirectional Reflectance Distribution Function (BRDF) models to capture highly directional scattering phenomena. Ultimately, UDMPM establishes a foundational blueprint for physics-aware neural rendering and 4D scene reconstruction.

\subsection{Limitations and Future Work}

While the proposed Unified Dual-Mask Physical Model (UDMPM) establishes a robust theoretical framework for non-Lambertian light field (LF) depth estimation, several limitations remain, which naturally dictate our future research directions from both signal processing and optical engineering perspectives.

First, the current MRI-inspired angular frequency analysis relies on a regular $9 \times 9$ angular sampling grid to perform 2D Discrete Fourier Transform (DFT). In scenarios with highly sparse or irregular angular sampling (e.g., $4 \times 4$ grids or wide-baseline multi-view setups), the angular frequency resolution degrades significantly. This causes spectral aliasing between high-frequency specular lobes and mid-frequency scattering components, potentially compromising the pixel-wise material classification and the accuracy of the medium mask ($M_{med}$). Furthermore, the current 2D-DFT assumes spatial stationarity within the local angular patch, which may introduce boundary artifacts at sharp material transitions. Future work will explore Non-Uniform Fast Fourier Transform (NUFFT) and graph signal processing techniques in the angular domain to achieve robust spectral parsing under arbitrary sampling topologies.

Second, although the dual-mask mechanism successfully unifies isotropic diffuse, specular, and scattering interactions via scalar surface roughness and 2D wavevector shifts, it inherently simplifies the high-dimensional Radiance Tensor Field (RTF). Complex material properties, such as anisotropic reflectance (e.g., brushed metals, hair) and sub-surface scattering (e.g., human skin, translucent plastics), cannot be fully parameterized by the current isotropic scalar and vector masks. To address this, we plan to extend the dual masks into high-order tensor representations and incorporate polarized light field cues. This extension will enable the physical decoding of anisotropic Bidirectional Reflectance Distribution Functions (BRDFs) and volumetric scattering, further generalizing the unified physical model to extremely complex real-world surfaces.

Finally, from an engineering perspective, the full inverse wavevector parsing and dual-mask modulation incur substantial computational and memory overhead when scaled to ultra-high-resolution LF data or dense video streams. Future efforts will focus on hardware-software co-design to alleviate this bottleneck. By integrating emerging computational imaging sensors, such as polarization-sensitive metasurfaces or event-based plenoptic cameras, we aim to directly capture high-frequency angular and photometric cues at the optical frontend. This would fundamentally bypass the spatial-angular resolution trade-off inherent in conventional microlens arrays, facilitating real-time, physically grounded 3D perception for next-generation autonomous navigation and robotic manipulation systems.

\bibliographystyle{IEEEtran}
\bibliography{references}

\end{document}
